% Options for packages loaded elsewhere
\PassOptionsToPackage{unicode}{hyperref}
\PassOptionsToPackage{hyphens}{url}
\documentclass[
]{article}
\usepackage{xcolor}
\usepackage{amsmath,amssymb}
\setcounter{secnumdepth}{-\maxdimen} % remove section numbering
\usepackage{iftex}
\ifPDFTeX
  \usepackage[T1]{fontenc}
  \usepackage[utf8]{inputenc}
  \usepackage{textcomp} % provide euro and other symbols
\else % if luatex or xetex
  \usepackage{unicode-math} % this also loads fontspec
  \defaultfontfeatures{Scale=MatchLowercase}
  \defaultfontfeatures[\rmfamily]{Ligatures=TeX,Scale=1}
\fi
\usepackage{lmodern}
\ifPDFTeX\else
  % xetex/luatex font selection
\fi
% Use upquote if available, for straight quotes in verbatim environments
\IfFileExists{upquote.sty}{\usepackage{upquote}}{}
\IfFileExists{microtype.sty}{% use microtype if available
  \usepackage[]{microtype}
  \UseMicrotypeSet[protrusion]{basicmath} % disable protrusion for tt fonts
}{}
\makeatletter
\@ifundefined{KOMAClassName}{% if non-KOMA class
  \IfFileExists{parskip.sty}{%
    \usepackage{parskip}
  }{% else
    \setlength{\parindent}{0pt}
    \setlength{\parskip}{6pt plus 2pt minus 1pt}}
}{% if KOMA class
  \KOMAoptions{parskip=half}}
\makeatother
\usepackage{longtable,booktabs,array}
\usepackage{caption}
\captionsetup[table]{skip=6pt}
\newcounter{none} % for unnumbered tables
\usepackage{calc} % for calculating minipage widths
% Correct order of tables after \paragraph or \subparagraph
\usepackage{etoolbox}
\makeatletter
\patchcmd\longtable{\par}{\if@noskipsec\mbox{}\fi\par}{}{}
\makeatother
\usepackage{graphicx}
\makeatletter
\newsavebox\pandoc@box
\newcommand*\pandocbounded[1]{% scales image to fit in text height/width
  \sbox\pandoc@box{#1}%
  \Gscale@div\@tempa{\textheight}{\dimexpr\ht\pandoc@box+\dp\pandoc@box\relax}%
  \Gscale@div\@tempb{\linewidth}{\wd\pandoc@box}%
  \ifdim\@tempb\p@<\@tempa\p@\let\@tempa\@tempb\fi% select the smaller of both
  \ifdim\@tempa\p@<\p@\scalebox{\@tempa}{\usebox\pandoc@box}%
  \else\usebox{\pandoc@box}%
  \fi%
}
% Set default figure placement to htbp
\def\fps@figure{htbp}
\makeatother
% definitions for citeproc citations
\NewDocumentCommand\citeproctext{}{}
\NewDocumentCommand\citeproc{mm}{%
  \begingroup\def\citeproctext{#2}\cite{#1}\endgroup}
\makeatletter
 % allow citations to break across lines
 \let\@cite@ofmt\@firstofone
 % avoid brackets around text for \cite:
 \def\@biblabel#1{}
 \def\@cite#1#2{{#1\if@tempswa , #2\fi}}
\makeatother
\newlength{\cslhangindent}
\setlength{\cslhangindent}{1.5em}
\newlength{\csllabelwidth}
\setlength{\csllabelwidth}{3em}
\newenvironment{CSLReferences}[2] % #1 hanging-indent, #2 entry-spacing
 {\begin{list}{}{%
  \setlength{\itemindent}{0pt}
  \setlength{\leftmargin}{0pt}
  \setlength{\parsep}{0pt}
  % turn on hanging indent if param 1 is 1
  \ifodd #1
   \setlength{\leftmargin}{\cslhangindent}
   \setlength{\itemindent}{-1\cslhangindent}
  \fi
  % set entry spacing
  \setlength{\itemsep}{#2\baselineskip}}}
 {\end{list}}
\usepackage{calc}
\newcommand{\CSLBlock}[1]{\hfill\break\parbox[t]{\linewidth}{\strut\ignorespaces#1\strut}}
\newcommand{\CSLLeftMargin}[1]{\parbox[t]{\csllabelwidth}{\strut#1\strut}}
\newcommand{\CSLRightInline}[1]{\parbox[t]{\linewidth - \csllabelwidth}{\strut#1\strut}}
\newcommand{\CSLIndent}[1]{\hspace{\cslhangindent}#1}
\setlength{\emergencystretch}{3em} % prevent overfull lines
\providecommand{\tightlist}{%
  \setlength{\itemsep}{0pt}\setlength{\parskip}{0pt}}
\makeatletter
\@ifpackageloaded{subcaption}{}{\usepackage{subcaption}}
\@ifpackageloaded{caption}{}{\usepackage{caption}}
\captionsetup[subfigure]{margin=0.5em}
\AtBeginDocument{%
\renewcommand*\figurename{Figure}
\renewcommand*\tablename{Table}
}
\AtBeginDocument{%
\renewcommand*\listfigurename{List of Figures}
\renewcommand*\listtablename{List of Tables}
}
\newcounter{pandoccrossref@subfigures@footnote@counter}
\newenvironment{pandoccrossrefsubfigures}{%
\setcounter{pandoccrossref@subfigures@footnote@counter}{0}
\begin{figure}\centering%
\gdef\global@pandoccrossref@subfigures@footnotes{}%
\DeclareRobustCommand{\footnote}[1]{\footnotemark%
\stepcounter{pandoccrossref@subfigures@footnote@counter}%
\ifx\global@pandoccrossref@subfigures@footnotes\empty%
\gdef\global@pandoccrossref@subfigures@footnotes{{##1}}%
\else%
\g@addto@macro\global@pandoccrossref@subfigures@footnotes{, {##1}}%
\fi}}%
{\end{figure}%
\addtocounter{footnote}{-\value{pandoccrossref@subfigures@footnote@counter}}
\@for\f:=\global@pandoccrossref@subfigures@footnotes\do{\stepcounter{footnote}\footnotetext{\f}}%
\gdef\global@pandoccrossref@subfigures@footnotes{}}
\@ifpackageloaded{float}{}{\usepackage{float}}
\floatstyle{ruled}
\@ifundefined{c@chapter}{\newfloat{codelisting}{h}{lop}}{\newfloat{codelisting}{h}{lop}[chapter]}
\floatname{codelisting}{Listing}
\newcommand*\listoflistings{\listof{codelisting}{List of Listings}}
\makeatother
\usepackage{bookmark}
\IfFileExists{xurl.sty}{\usepackage{xurl}}{} % add URL line breaks if available
\urlstyle{same}
% fallback for those not using the hyperref driver hyperxmp:
\makeatletter
\@ifundefined{xmpquote}{\newcommand{\xmpquote}[1]{#1}}{}
\makeatother
\hypersetup{
  hidelinks,
  pdfcreator={LaTeX via pandoc}}

\author{}
\date{}

\begin{document}

\section*{Correction to 10.1029/2026JB033781}\label{sec:correction}
\addcontentsline{toc}{section}{Correction to 10.1029/2026JB033781}

The kinetic prefactor \(Z\) in the interface-controlled olivine
\(\Leftrightarrow\) wadsleyite rate law (Equation 13 in
10.1029/2026JB033781) was reported with units of K\(^{-1}\) s\(^{-1}\):

\begin{equation}\protect\phantomsection\label{eq:reaction-rate}{\dot{X} = Z\, T\, \exp\!\left(-\frac{H^{\ast} + P V^{\ast}}{R T}\right) \left(1 - \exp\!\left[-\frac{\Delta G}{R T}\right]\right) \left(1 - X\right)}\end{equation}

Because our ASPECT simulations were set to use time in years, \(Z\) was
inadvertently scaled by a factor of 3.15e7. This was user error in
parameter passing, not a bug in ASPECT. The published range \(Z\) =
3.0e0--7.0e7 K\(^{-1}\) s\(^{-1}\) effectively ran as \(Z^{\prime}\) =
9.5e-8--2.2e0 K\(^{-1}\) s\(^{-1}\).

The \(Z\) values have been corrected everywhere they appear in
10.1029/2026JB033781 and the Supplementary Information. We have edited
two passages in Sections 2.2.3.2 and 4.1 to contextualize the correction
in terms of the experimental constraints of Hosoya et al. (2005). The
accompanying dataset on OSF (Kerswell, 2026a) is read-only and has not
changed. The latest code release on Zenodo (Kerswell, 2026b) reproduces
the corrected manuscript directly from the archived OSF dataset (by
dividing \(Z\) by 3.15e7). All else remains unchanged.

Every result comes directly from the simulation output, so all figures,
regime thresholds, and conclusions stand as published in
10.1029/2026JB033781. However, they reflect slower, drier kinetics than
the original text implied.

Hosoya et al. (2005) reports \(\ln \Gamma\) = -18.0 \(\pm\) 3.8 (assumed
1\(\sigma\)). With \(Z = \frac{6.67}{d}\, \Gamma\, C_\mathrm{OH}^n\) and
\(n\) = 3.2, \(C_\mathrm{OH}\) = 50--5000 ppm, \(d\) = 1--10 mm, the
permitted \(Z\) range across \(\pm 2\sigma\) runs from 1.4e-3 K\(^{-1}\)
s\(^{-1}\) (10-mm grains, 50 ppm OH at \(-2\sigma\)) to 1.4e11
K\(^{-1}\) s\(^{-1}\) (1-mm grains, 5000 ppm OH at \(+2\sigma\)).
Extending to 1 ppm OH widens the full \(\pm 2\sigma\) range to
5.1e-09--1.4e11 K\(^{-1}\) s\(^{-1}\). The published (3.0e0--7.0e7
K\(^{-1}\) s\(^{-1}\)) and effective (9.5e-8--2.2e0 K\(^{-1}\)
s\(^{-1}\)) ranges both fall inside the extended dry conditions (Figure
\ref{fig:erratum-parameter-ranges}). Our simulations therefore
effectively sampled conditions from nominally dry, coarse-grains to
relatively dry (hundreds of ppm OH), fine-grains. This is consistent
with Ishii \& Ohtani (2021), who showed that olivine and wadsleyite can
remain kinetically dry even in hydrated slabs.

Under ambient mantle conditions at the top of the domain (\(P\) = 10
GPa, \(T\) = 1706 K, \(H^{\ast}\) = 274 kJ mol\(^{-1}\), \(V^{\ast}\) =
3.0e-6 m\(^3\) mol\(^{-1}\)), published \(Z\) values give transformation
times of milliseconds to days, and effective \(Z^{\prime}\) values give
0.02 yr to 0.4 Ma. Our simulations therefore retain metastable olivine
at warmer slab temperatures (1350--1400 K at 15 GPa, or 500 K colder
than the ambient mantle) due to sluggish kinetics. This contrasts with
earlier work that used different kinetics and attributed metastable
olivine to much colder slab temperatures (773--923 K at 15 GPa, Rubie \&
Ross, 1994; Däßler et al., 1996; Devaux et al., 1997; Hosoya et al.,
2005).

\floatplacement{figure}{H}

\begin{figure}
\centering
\includegraphics[width=1\linewidth,height=\textheight,keepaspectratio,alt={Published
Z range (red) and effective Z\^{}\{\textbackslash prime\} range passed
to ASPECT (blue). The thin solid box marks the experimental conditions
of Hosoya et al. (2005); the dashed box marks extended dry conditions
(down to 1 ppm OH). Dashed contours are calculated Z values. Thick solid
contours mark the permitted minimum (left panel, 5.1e-09 K\^{}\{-1\}
s\^{}\{-1\}) and maximum (right panel, 1.4e11 K\^{}\{-1\} s\^{}\{-1\}) Z
values for the extended dry conditions. The blue and red ranges overlap
at the 95\% confidence interval (\textbackslash pm 2\textbackslash sigma
\textbackslash ln \textbackslash Gamma). The arrow (middle panel) marks
the accidental 3.15e7 scaling in Z.}]{../figs/erratum-kinetics.png}
\caption{Published \(Z\) range (red) and effective \(Z^{\prime}\) range
passed to ASPECT (blue). The thin solid box marks the experimental
conditions of Hosoya et al. (2005); the dashed box marks extended dry
conditions (down to 1 ppm OH). Dashed contours are calculated \(Z\)
values. Thick solid contours mark the permitted minimum (left panel,
5.1e-09 K\(^{-1}\) s\(^{-1}\)) and maximum (right panel, 1.4e11
K\(^{-1}\) s\(^{-1}\)) \(Z\) values for the extended dry conditions. The
blue and red ranges overlap at the 95\% confidence interval
(\(\pm 2\sigma \ln \Gamma\)). The arrow (middle panel) marks the
accidental 3.15e7 scaling in \(Z\).}\label{fig:erratum-parameter-ranges}
\end{figure}

\section*{References}\label{sec:references}
\addcontentsline{toc}{section}{References}

\protect\phantomsection\label{refs}
\begin{CSLReferences}{1}{0}
\bibitem[\citeproctext]{ref-dassler1996b}
Däßler, R., Yuen, D., Karato, S., \& Riedel, M. (1996). Two-dimensional
thermo-kinetic model for the olivine-spinel phase transition in
subducting slabs. \emph{Physics of the Earth and Planetary Interiors},
\emph{94}(3-4), 217--239.

\bibitem[\citeproctext]{ref-devaux1997}
Devaux, J., Schubert, G., \& Anderson, C. (1997). Formation of a
metastable olivine wedge in a descending slab. \emph{Journal of
Geophysical Research: Solid Earth}, \emph{102}(B11), 24627--24637.

\bibitem[\citeproctext]{ref-hosoya2005}
Hosoya, T., Kubo, T., Ohtani, E., Sano, A., \& Funakoshi, K. (2005).
Water controls the fields of metastable olivine in cold subducting
slabs. \emph{Geophysical Research Letters}, \emph{32}(17).

\bibitem[\citeproctext]{ref-ishii2021}
Ishii, T., \& Ohtani, E. (2021). Dry metastable olivine and slab
deformation in a wet subducting slab. \emph{Nature Geoscience},
\emph{14}(7), 526--530.

\bibitem[\citeproctext]{ref-kerswell2026a}
Kerswell, B. (2026a, April).
{buchanankerswell/kerswell\_et\_al\_410\_kinetics: Publication Release
{[}Dataset{]}} (Version v.1.0.0). OSF.
\url{https://doi.org/10.17605/OSF.IO/9PHWC}

\bibitem[\citeproctext]{ref-kerswell2026c}
Kerswell, B. (2026b, October).
{buchanankerswell/kerswell\_et\_al\_410\_kinetics: Publication Release
{[}Software{]}} (Version v.1.1.0). Zenodo.
\url{https://doi.org/10.5281/zenodo.19661753}

\bibitem[\citeproctext]{ref-rubie1994}
Rubie, D., \& Ross, C. (1994). Kinetics of the olivine-spinel
transformation in subducting lithosphere: Experimental constraints and
implications for deep slab processes. \emph{Physics of the Earth and
Planetary Interiors}, \emph{86}(1-3), 223--243.

\end{CSLReferences}

\end{document}
